\documentclass[11pt,a4paper,UTF8]{ctexart}
\usepackage{geometry}
\geometry{left=18mm,right=18mm,top=24mm,bottom=24mm,headheight=22pt,footskip=12mm}
\usepackage{amsmath,amssymb,mathtools,bm,esint,extarrows,centernot}
\usepackage{xcolor}
\usepackage{tcolorbox}
\tcbuselibrary{skins,breakable}
\usepackage{fancyhdr}
\usepackage{titlesec}
\usepackage{hyperref}
\usepackage{tikz}
\usepackage{booktabs}
\usepackage{tabularx}

% --- Color Palette (Purple & DeepPink Academic Theme) ---
\definecolor{DeepPurple}{HTML}{4C1D95}
\definecolor{TitlePurple}{HTML}{581C87}
\definecolor{SubPurple}{HTML}{6B21A8}
\definecolor{DeepPink}{HTML}{FF1493}
\definecolor{LightPinkBg}{HTML}{FFF5F8}
\definecolor{LightPurpleBg}{HTML}{F8F5FF}
\definecolor{GoldAccent}{HTML}{D97706}
\definecolor{DarkText}{HTML}{1F2937}

\hypersetup{
    colorlinks=true,
    linkcolor=TitlePurple,
    citecolor=DeepPink,
    urlcolor=SubPurple,
    pdfborder={0 0 0}
}

% --- Typography & Spacing ---
\linespread{1.3}
\setlength{\parskip}{0.35em plus 0.1em minus 0.1em}
\setlength{\parindent}{0pt}

% --- Section Titles Styling ---
\titleformat{\section}
  {\Large\bfseries\color{TitlePurple}}{\thesection}{1em}{}
  [\vspace{1mm}{\color{TitlePurple!30}\hrule height 1pt}]
\titleformat{\subsection}
  {\large\bfseries\color{SubPurple}}{\thesubsection}{0.8em}{}
\titleformat{\subsubsection}
  {\normalsize\bfseries\color{DeepPurple}}{\thesubsubsection}{0.6em}{}

% --- tcolorbox Environments ---
% 1. Theorem / Formula / Method / Analysis Box (Deep Purple)
\newtcolorbox{thmbox}[1][]{
    enhanced,
    breakable,
    colback=LightPurpleBg,
    colframe=TitlePurple,
    arc=3pt,
    boxrule=0pt,
    leftrule=3.5pt,
    left=10pt,
    right=10pt,
    top=8pt,
    bottom=8pt,
    title=#1,
    coltitle=white,
    colbacktitle=TitlePurple,
    fonttitle=\bfseries\small,
    attach boxed title to top left={yshift=-2mm, xshift=4mm},
    boxed title style={boxrule=0pt, arc=2pt}
}

% 2. Solution / Formula / Derivation Box (DeepPink)
\newtcolorbox{solbox}[1][]{
    enhanced,
    breakable,
    colback=LightPinkBg,
    colframe=DeepPink,
    arc=3pt,
    boxrule=0pt,
    leftrule=3.5pt,
    left=10pt,
    right=10pt,
    top=8pt,
    bottom=8pt,
    title=#1,
    coltitle=white,
    colbacktitle=DeepPink,
    fonttitle=\bfseries\small,
    attach boxed title to top left={yshift=-2mm, xshift=4mm},
    boxed title style={boxrule=0pt, arc=2pt}
}

% 3. Learning Goals / Summary Box
\newtcolorbox{targetbox}[1][]{
    enhanced,
    breakable,
    colback=LightPurpleBg,
    colframe=DeepPurple,
    arc=3pt,
    boxrule=1pt,
    left=10pt,
    right=10pt,
    top=8pt,
    bottom=8pt,
    title=#1,
    coltitle=white,
    colbacktitle=DeepPurple,
    fonttitle=\bfseries\small,
    attach boxed title to top left={yshift=-2mm, xshift=4mm},
    boxed title style={boxrule=0pt, arc=2pt}
}

\begin{document}

% ------------------- 讲义封面 (Cover Page) -------------------
\begin{titlepage}
\thispagestyle{empty}
\begin{tikzpicture}[remember picture,overlay]
    \fill[DeepPurple] (current page.north west) rectangle ([yshift=-7cm]current page.north east);
    \draw[DeepPink, line width=3.5pt] ([yshift=-7cm]current page.north west) -- ([yshift=-7cm]current page.north east);
    \draw[GoldAccent, line width=1pt] ([yshift=-7.12cm]current page.north west) -- ([yshift=-7.12cm]current page.north east);
    
    \node[anchor=north east, opacity=0.12, text=white] at ([xshift=-1.5cm, yshift=-0.5cm]current page.north east) {
        \fontsize{70}{70}\selectfont\bfseries 2027
    };
\end{tikzpicture}

\vspace*{0.8cm}
\begin{center}
    {\color{white}\fontsize{26}{32}\selectfont\bfseries 2027 考研数学(二) 核心讲义}\\[0.6cm]
    {\color{white}\fontsize{13}{18}\selectfont\bfseries 全国硕士研究生招生考试 · 统考数学(二) 核心精讲全书}\\[1.5cm]
\end{center}

\vspace{3.5cm}
\begin{center}
    {\color{GoldAccent}\large\bfseries $\blacktriangleright$ 基础强化 · 核心题解 · 考点深水区 $\blacktriangleleft$}\\[0.8cm]
    {\color{TitlePurple}\fontsize{24}{28}\selectfont\bfseries 第四章 一元函数微分学的计算}\\[0.6cm]
    {\color{DeepPink}\large\bfseries —— 体系化理论精讲与公式全景 ——}\\[1.5cm]
\end{center}

\vfill

\begin{center}
\begin{tcolorbox}[
    colback=white,
    colframe=DeepPurple,
    width=0.88\textwidth,
    arc=4pt,
    boxrule=1.2pt,
    halign=center,
    drop shadow=gray!30
]
    \vspace{3mm}
    {\bfseries\large 编著团队：EscoffierZhou}\\[2.5mm]
    {\bfseries\color{TitlePurple} 目标院校：中国科学院大学 (UCAS) 计算机专硕 (22408)}\\[2.5mm]
    {\small\color{gray} 备考铁律：手算真题数字 · 错题白纸重做 · 408 客观题为王}\\[2.5mm]
    {\small 编制版本：2027 届首发版 · \today}
    \vspace{2mm}
\end{tcolorbox}
\end{center}
\vspace{0.8cm}
\end{titlepage}

% ------------------- 目录与页面主体配置 -------------------
\pagestyle{fancy}
\fancyhf{}
\fancyhead[L]{\small\color{gray}2027 考研数学(二) · 核心讲义系列}
\fancyhead[R]{\small\color{TitlePurple}\nouppercase{\leftmark}}
\fancyfoot[C]{\small\color{gray}— 第 \thepage\ 页 —}
\fancyfoot[R]{\small\color{gray}UCAS 22408}
\renewcommand{\headrulewidth}{0.5pt}

\pagenumbering{Roman}
\tableofcontents
\newpage
\pagenumbering{arabic}


\begin{targetbox}[本章复习目标与核心知识图谱]
\begin{itemize}
  \item 目的[1]:掌握四则运算和复合函数求导/初等函数的求数公式
\end{itemize}
\end{targetbox}

\begin{targetbox}[本章复习目标与核心知识图谱]
\begin{itemize}
  \item 目的[2]:掌握微分四则运算/一阶微分形式的不变性
\end{itemize}
\end{targetbox}

\begin{targetbox}[本章复习目标与核心知识图谱]
\begin{itemize}
  \item 目的[3]:求分段函数的导数/隐函数+参数方程确定的函数/反函数的导数
\end{itemize}
\end{targetbox}


\section{1.基本初等函数求导公式与四则运算}

\begin{thmbox}[原理[1]: 基本求导函数]

$(x^\alpha)' = \alpha x^{\alpha-1}$	$(a^x)' = a^x \ln a$	$(e^x)' = e^x$

$(\log_a x)' = \frac{1}{x \ln a}$	$(\ln\vert{}x\vert{})' = \frac{1}{x} \ (x \neq 0)$

\textit{}*

\end{thmbox}
\begin{thmbox}[原理[2]: 三角函数/反三角函数求导函数]

($secx=\frac{1}{cosx}$,$cscx = \frac{1}{sinx}$,$cotx=\frac{1}{tanx}$)

$(\sin x)' = \cos x$	$(\cos x)' = -\sin x$	$(\tan x)' = \sec^2 x$

$(\cot x)' = -\csc^2 x$    [$(\frac{1}{tanx})'=-\frac{1}{sin^2x}$]     

$(\sec x)' = \sec x \tan x$  [$(\frac{1}{cosx})'=\frac{tanx}{cosx}$] 

$(\csc x)' = -\csc x \cot x$ [$(\frac{1}{sinx})'=-\frac{1}{sinxtanx}$]

$(\arcsin x)' = \frac{1}{\sqrt{1-x^2}}$ $(\arccos x)' = -\frac{1}{\sqrt{1-x^2}}$   

$(\arctan x)' = \frac{1}{1+x^2}$  $(\operatorname{arccot} x)' = -\frac{1}{1+x^2}$	

\textbf{}

\end{thmbox}
\begin{thmbox}[原理[3]: 双曲反函数求导函数]

$[\ln(x + \sqrt{x^2+a^2})]' = \frac{1}{\sqrt{x^2+a^2}}$（特别地，$a=1$时为$\frac{1}{\sqrt{x^2+1}}$）

$[\ln(x + \sqrt{x^2-1})]' = \frac{1}{\sqrt{x^2-1}}$

\textbf{}

\end{thmbox}
\begin{thmbox}[原理[4]: 导数的四则运算]

设$u=u(x), v=v(x)$均可导：

和:$[u(x) + v(x)]' = u'(x) + v'(x)$

差:$[u(x) - v(x)]' = u'(x) - v'(x)$

积:$[u(x)v(x)]' = u'(x)v(x) + u(x)v'(x)$

 $[uvw]' = u'vw + uv'w + uvw'$

商:$\left[\frac{u(x)}{v(x)}\right]' = \frac{u'(x)v(x) - u(x)v'(x)}{[v(x)]^2} \ (v(x) \neq 0)$

\textit{}*

\end{thmbox}

\begin{equation*}
\{f[g(x)]\}' = f'[g(x)] \cdot g'(x) \quad \left(\text{即 } \frac{\mathrm{d}y}{\mathrm{d}x} = \frac{\mathrm{d}y}{\mathrm{d}u} \cdot \frac{\mathrm{d}u}{\mathrm{d}x}\right)
\end{equation*}
\begin{thmbox}[原理[5]: 复合函数求导]

原理:设$y=f(u)$在点$u$可导,$u=g(x)$在点$x$可导，

则复合函数$y=f[g(x)]$在点$x$可导,且:


\textit{}*

\end{thmbox}
\begin{thmbox}[原理[6]: 分段函数求导]

设$f(x) = \begin{cases} f_1(x), & x \ge x_0 \\ \\ f_2(x), & x < x_0 \end{cases}$

对非分段点($x \neq x_0$):直接求导

对于分段点($x = x_0$):[1]用导数定义求左、右导数[2]如果在$x_0$连续,则$f'_+(x_0) = \lim_{x \to x_0^+} f'(x)$

  
\begin{equation*}
f'_+(x_0) = \lim_{x \to x_0^+} \frac{f_1(x) - f(x_0)}{x - x_0}, \quad f'_-(x_0) = \lim_{x \to x_0^-} \frac{f_2(x) - f(x_0)}{x - x_0}
\end{equation*}


  若$f'_+(x_0) = f'_-(x_0)$,则导数存在,且$f'(x_0)$等于该极限；否则不可导

\textit{}*

\end{thmbox}
\begin{thmbox}[原理[7]: 反函数求导]

设$y=f(x)$单调、可导且$f'(x) \neq 0$,其反函数为$x=\varphi(y)$：

一阶导数：$\frac{\mathrm{d}x}{\mathrm{d}y} = \frac{1}{\frac{\mathrm{d}y}{\mathrm{d}x}}$，即$\varphi'(y) = \frac{1}{f'(x)}$

二阶导数：
\begin{equation*}
\varphi''(y) = \frac{\mathrm{d}}{\mathrm{d}y}\left(\frac{1}{f'(x)}\right) = \frac{\mathrm{d}}{\mathrm{d}x}\left(\frac{1}{f'(x)}\right) \cdot \frac{\mathrm{d}x}{\mathrm{d}y} = -\frac{f''(x)}{[f'(x)]^2} \cdot \frac{1}{f'(x)} = -\frac{f''(x)}{[f'(x)]^3}
\end{equation*}


\textit{}*

\end{thmbox}
\begin{thmbox}[原理[8]: 隐函数求导]

操作[1]:由方程$F(x, y) = 0$确定的可导函数$y = y(x)$：

操作[2]:方程两边直接对$x$求导

  注意将$y$视为$x$的函数（如$(y^2)' = 2y \cdot y'$，$(xy)' = y + x y'$）

操作[3]:解出含$y'$的代数方程即可求得$y'$

\textit{}*

\end{thmbox}
\begin{thmbox}[原理[9]: 参数函数求导]

设曲线由参数方程$\begin{cases} x = \varphi(t) \\ y = \psi(t) \end{cases}$确定，其中$\varphi(t), \psi(t)$可导且$\varphi'(t) \neq 0$：

\textbf{一阶导数}：

\end{thmbox}

\begin{equation*}
\frac{\mathrm{d}y}{\mathrm{d}x} = \frac{\frac{\mathrm{d}y}{\mathrm{d}t}}{\frac{\mathrm{d}x}{\mathrm{d}t}} = \frac{\psi'(t)}{\varphi'(t)}
\end{equation*}

\textbf{二阶导数}：


\begin{equation*}
\frac{\mathrm{d}^2y}{\mathrm{d}x^2} = \frac{\mathrm{d}}{\mathrm{d}x}\left(\frac{\mathrm{d}y}{\mathrm{d}x}\right) = \frac{\frac{\mathrm{d}}{\mathrm{d}t}\left(\frac{\mathrm{d}y}{\mathrm{d}x}\right)}{\frac{\mathrm{d}x}{\mathrm{d}t}} = \frac{\psi''(t)\varphi'(t) - \psi'(t)\varphi''(t)}{[\varphi'(t)]^3}
\end{equation*}

\textit{}*

\begin{thmbox}[原理[10]: 对数/幂指求导法]

(适用于多项乘除/开方乘方/幂指函数$y = u(x)^{v(x)}$)

\textbf{操作[1]:对数求导法}

两边取自然对数$\ln y = v(x)\ln u(x)$，两边对$x$求导得$\frac{y'}{y} = [v(x)\ln u(x)]'$，从而$y' = y \cdot [v(x)\ln u(x)]'$

\textbf{操作[2]化为复合指数求导:}

写成$u(x)^{v(x)} = e^{v(x)\ln u(x)}$后直接用复合函数链式法则

\end{thmbox}

\section{2.微分运算与一阶微分形式不变性}

\begin{thmbox}[原理[6]: 微分四则运算]

$\mathrm{d}(u + v) = \mathrm{d}u + \mathrm{d}v$

$\mathrm{d}(u - v) = \mathrm{d}u - \mathrm{d}v$

$\mathrm{d}(uv) = v\,\mathrm{d}u + u\,\mathrm{d}v$

$\mathrm{d}\left(\frac{u}{v}\right) = \frac{v\,\mathrm{d}u - u\,\mathrm{d}v}{v^2} \ (v \neq 0)$

\end{thmbox}

\begin{equation*}
\mathrm{d}y = f'(u)\,\mathrm{d}u
\end{equation*}
\begin{thmbox}[原理[7]: 一阶微分形式不变性]

无论是$x$作为自变量,还是$u=g(x)$作为中间变量,$y=f(u)$的微分形式始终保持：


  一阶微分形式不变性是凑微分法与链式法则的核心基础

  微分运算:不断使用基本公式:$dy = f'(y)du = f'(y)f'(u)dx$,使得化简为dx

\end{thmbox}
\begin{thmbox}[原理[8]: 凑微分(反向换元)]

比如$\int\frac{2x}{x^2+1}dx$,其中$d(x^2+1)=2xdx$,所以$I=\frac{d(x^2+1)}{x^2+1}=ln(x^2+1)+C$

\end{thmbox}

\end{document}
